% Folder 79, batch 08 — archivists' pages 141–160.
% Pass: Opus 5.5 (claude-opus-5-5), 2026-10-03 — first pass, unchecked against the pages by a human.
%
% His own pagination runs 72 to 81 on the rectos (p. 141 = 72, p. 143 = 73,
% p. 145 = 74, p. 147 = 75, p. 149 = 76, p. 151 = 77, p. 153 = 78, p. 155 = 79,
% p. 157 = 80, p. 159 = 81). Page 158 is half blank; page 159 opens § IX.
\documentclass[11pt,a4paper]{article}
\input{../preamble/grothendieck}

\folder{79}
\batch{8}
\pages{141}{160}
\foldertitle{2-polyèdres réguliers triangulés, et modules libres de rangs 2 sur les anneaux $\Lambda=\mathbb{Z}/n\mathbb{Z}$ et $\Lambda=\mathbb{Z}$ : notes manuscrites (s.d.).}
\dating{[à partir de 1980]}
\watermark{Édition de démonstration}

\begin{document}

\page{141}\note{p. 72 de l'auteur. La page s'ouvre au milieu d'une phrase commencée avant ce lot.}
pour son propre mérite. La question bien \uncertain{sûr} ne se pose que \uncertain{si} $Z'_n \neq \{1\}$, i.e. si $n \notin \{2,3,4,6\}$, le premier cas est donc celui où
\[
(37) \qquad n = 5, \quad \Lambda'_n \simeq \mathbb{F}_5^{*}/\pm 1 \simeq \mathbb{Z}/2\mathbb{Z},
\]
considérons une structure de 2-polyèdre comme triangulé (= faces d'ordre 3), dont les sommets sont d'ordre 5. On sait qu'il n'y a que l'icosaèdre (structure orientable : 12 sommets) et l'icosaèdre \uncertain{géminé} (structure \uncertain{non orientée}). L'opération envisagée ici (avec l'élément non trivial de $\Lambda'_5$) est celle qui fait passer d'un icosaèdre à l'icosaèdre associé — l'ens.\ $\pi$ (dans le cas orienté) étant la structure de \emph{bi-icosaèdre}. L'icosaèdre peut être décrit : à l'aide d'un \uncertain{module} $M$ de \uncertain{rang} 2 sur $\mathbb{F}_5$, comme $(S_M, F_M)$, et satisfait donc à l'axiomatique du présent §. Mais l'icosaèdre \uncertain{géminé} ? Notons que le \uncertain{commutant} de $G_0$ dans ce cas se réduit à l'identité, donc $Z = 1$, donc $Z \to \Lambda'^{*}_n$ \ill{} \uncertain{surjectif} — \ill{} \struck{\ill{}} \ill{} l'opération $A_0 \to A_0^{\vee}$, dans ce cas-là ! Cela

\page{142}
montre bien que dans certains cas, la validité de l'opération $A_0 \mapsto A_0^{\vee}$ dépend de celle de l'axiomatique T~1 à T~5 \uncertain{discutée} \ill{} présent.

Le cas suivant serait \struck{\ill{}}
\[
n = 7, \qquad Z'_n \simeq \mathbb{Z}/3\mathbb{Z} \qquad \bigl((\mathbb{Z}/7\mathbb{Z})^{*} \simeq \mathbb{Z}/6\mathbb{Z}\bigr)
\]
engendré par n'importe quel $\xi \in \mathbb{F}_7^{*} \setminus \{1,-1\} = \mathbb{F}_7 \setminus \{0,1,-1\}$, p.ex.\ $\nu = 2 \bmod 7$, dont l'inverse est $\nu' = 4 \bmod 7$. Est-il vrai que pour tout 2-polyèdre $(S, A_0, F_0)$ \struck{\ill{}} triangulé, à sommets d'ordre 7, le \uncertain{quotient} $\bigl(S, A_0^{(2)} \subset \mathcal{P}_2(S)\bigr)$ correspondant \uncertain{soit} encore un tel 2-polyèdre ? Même question pour $n = 8$, où $Z'_8 \simeq \mathbb{Z}/2\mathbb{Z}$, ou $n = 10$ \uncertain{où} $Z'_{10} \simeq \mathbb{Z}/2\mathbb{Z}$ (pour $n = 9$ on a $Z'_9 \simeq \mathbb{Z}/4\mathbb{Z}$ — en fait pour
\[
n \notin \{\underbrace{2,3,4,6}_{Z'_n = 1}, \underbrace{5,8,10,12}_{Z'_n \simeq \mathbb{Z}/2\mathbb{Z}}\} \quad \text{i.e.} \quad n \geqslant 13 \ \text{ou}\ n \in \{7, 9, 11\},
\]
\note{sous « 7 » : $Z'_n \simeq \mathbb{Z}/3\mathbb{Z}$ ; un trait relie « 9, 11 » à $Z'_n \in \mathbb{Z}/4\mathbb{Z}$, sic ; le « 11 » est récrit sur un autre chiffre}
$Z'_n$ d'ordre $\geqslant 3$ — ordre 3 pour $n = 7$ et 14, ordre 4 pour $n = 9$ et 18 (et pour $n = 16, 20$), ordre 5 pour $n = 11$ et 22, — ordre 6 pour $n = 13$ et 26 — (ordre 12 pour $n = 15$ ou 30) \struck{\ill{}} ordre 9 pour \struck{\ill{}} 19 et 38, 8 pour $n = 17$ ou 34,\note{« (ordre 12 pour $n = 15$ ou 30) » est rattaché par une accolade aux ordres de la ligne ; pour $n = 15$, $(\mathbb{Z}/15\mathbb{Z})^{*}/\pm 1$ est d'ordre 4 : la lecture ou le calcul est douteux}
\marginal{8 \uncertain{(?)}, en tête de la dernière ligne}

Mais poursuivons dans le cadre de l'axiomatique T~1 …, en \uncertain{trouvant}
\[
(38) \qquad (\mathrm{T}\,7) \quad \text{L'homomorphisme } Z \to \Lambda'^{*}_n \text{ est surjectif, i.e. un isomorphisme.}
\]

\page{143}\note{p. 73 de l'auteur.}
\marginal{\ill{} — note de huit lignes environ, écrite en diagonale dans le coin supérieur gauche ; seuls quelques mots se laissent deviner (« il vaut mieux », « \uncertain{l'axiomatique} »)}
(39) (T~8) Le 2-polyèdre $P_0$ est régulier, et orientable

— ce qui fait qu'on doit pouvoir traduire l'axiomatique en termes de théorie des groupes, savoir en termes des propriétés d'un groupe $G_0^{+}$ qui est engendré par des générateurs $\rho_0, \rho_1, \rho_\infty$ satisfaisant
\[
(40) \qquad \boxed{\ \rho_0^{n} = \rho_1^{2} = \rho_\infty^{3} = 1\ }
\]
et où $\rho_0$ est d'ordre exactement $n$, sous-groupe \uncertain{d'indice} 2 d'un groupe $G_0$ engendré par des générateurs $\tau_0, \tau_1, \tau_\infty$ reliés entre eux et aux précédents par
\[
(41) \qquad \begin{cases} \tau_0^{2} = \tau_1^{2} = \tau_\infty^{2} = 1 \\ \rho_0 = \tau_\infty \tau_1, \quad \rho_1 = \tau_0 \tau_\infty, \quad \rho_\infty = \tau_1 \tau_0 . \end{cases}
\]

Je vais d'abord \struck{introduire} \add{introduire le} \struck{Un \ill{} isomorphisme} \add{$\Lambda_n$-torseur}
\[
(42) \qquad \vec{\pi} = \coprod_{i \in \pi} \mu(i), \qquad \text{où } \mu(i) = \text{ens.\ des deux orientations de } P_i = (S, A_i, F_i)
\]
on fait opérer dessus \struck{$(\mathbb{Z}/n\mathbb{Z})^{*} \simeq \Lambda$}
\[
(43) \qquad Z_n = \Lambda_n^{*} = (\mathbb{Z}/n\mathbb{Z})^{*}
\]
en notant que pour tout $s \in S$, $i \in \pi$, \struck{\ill{}} $\dot\lambda \in Z' \simeq Z'_n$, on a sur $\mathrm{Et}_i(s) = \mathrm{Et}_{\dot\lambda . i}(\uncertain{s})$, deux structures polygonales $\alpha_i(s)$, $\alpha_{\dot\lambda . i}(s)$, reliées par
\[
(44) \qquad \alpha_{\dot\lambda . i}(s) = \alpha_i(s)^{\dot\lambda},
\]

\page{144}
ce qui, par la définition de l'opération de $Z'_n$ sur \uncertain{l'ensemble} des structures polygonales sur un ensemble $E$, \struck{\ill{}} implique \uncertain{qu'on a} un isomorphisme canonique
\[
(43) \qquad p^{-1}(\dot\lambda) \wedge \mu(i) \wedge \mu(\dot\lambda i) \simeq \mathbb{1}
\]
\marginal{sous la formule : torseurs \uncertain{triviaux} sous $\pm 1$}
où
\[
p : Z_n = \Lambda_n^{*} \longrightarrow Z'_n = \Lambda_n^{*}/\pm 1
\]
est l'homomorphisme canonique, \struck{\ill{}} de degré 2 (NB on suppose $n \neq 2$ donc $1 \neq -1$ dans $\Lambda_n$). On a en fait pour
\[
\lambda \in p^{-1}(\dot\lambda), \quad \omega \in \mu(i), \quad \omega' \in \mu(i' = \dot\lambda . i)
\]
\[
(44) \qquad \lambda \wedge \omega \wedge \omega' = 1 \iff \omega'^{\,\uncertain{\lambda}}_{E} = \omega_E^{\lambda}
\]
\note{le numéro et l'exposant de $\omega'_E$ sont surchargés ; on attend $\omega'_E = \omega_E^{\lambda}$, comme en (46)}
où $\omega_E, \omega'_E$ sont \struck{\ill{}} les permutations circulaires de l'ensemble \uncertain{étoilé} $E$ qui correspondent aux orientations $\omega, \omega'$.

A priori l'isomorphisme (43) dépend du choix du sommet $s$. Mais par transport de structure, il est le même pour deux sommets $s, s'$ tels que l'on ait
\[
s' = g . s, \qquad \text{avec } g \in G_0^{+},
\]
à condition de supposer que $g$ opère sur $\mu(i)$ et sur $\mu(\dot\lambda i)$ \uncertain{avec la même signature}. Si on admet que c'est \uncertain{toujours} le cas (cf plus bas), il suffit de prendre $g \in G_0^{+}$, \struck{\ill{}} \add{\ill{} orientation commune}
\[
(45) \qquad G_0 \xrightarrow{\ \mathrm{sg}\ } \pm 1 ,
\]

\page{145}\note{p. 74 de l'auteur.}
\marginal{en diagonale, à gauche : Il est clair que $\forall\, \omega, \omega' \in \vec{\pi}$, $\exists !\, \lambda \in \Lambda_n^{*}$, $\omega' = \lambda . \omega$}
et comme $G_0^{+}$ est transitif sur $S$, il s'ensuit bien que (43) est indépendant du choix de $s \in S$. [Si on ne supposait pas les $P_i$ \emph{réguliers}, il y aurait lieu soit de mettre un axiome \uncertain{d'indépendance}, soit de la justifier par voie combinatoire, en termes des autres propriétés …] Ceci dit, si $(\lambda, \omega, \omega') \in \Lambda_n \times \mu(i) \times \mu(\dot\lambda . i)$ se correspondent par (44) (pour un $s$, donc pour tout $s$), on posera par définition
\[
(46) \qquad \begin{cases} \omega' = \lambda . \omega \qquad (\omega \in \mu(i),\ \omega' \in \mu(i') = \mu(\dot\lambda . i)) \\ \iff \omega'_E = \omega_E^{\lambda} \ \text{pour tout étoilé (ou un étoilé)} \qquad \lambda \in \Lambda_n^{*} ). \end{cases}
\]
Soit de même $\lambda' \in \Lambda_n^{*}$, $\omega'' = \lambda' . \omega'$, \uncertain{$\in \mu(i'')$}, $i'' = (\dot\lambda' \dot\lambda) i = \dot\lambda' i'$, considérons
\[
\omega''_E = (\omega'_E)^{\lambda'} \quad \text{donc} \quad \omega''_E = \bigl((\omega_E)^{\lambda}\bigr)^{\lambda'} = (\omega_E)^{\lambda\lambda'}
\]
on a donc aussi, par définition de $(\lambda\lambda') . \omega$,
\[
\omega'' = (\lambda\lambda') . \omega
\]
\struck{\ill{}} d'où la relation d'associativité
\[
(47) \qquad \lambda'(\lambda \omega) = (\lambda'\lambda) . \omega ,
\]
donc la loi (46) définit bien une opération du groupe $\Lambda_n^{*}$ sur $\vec{\pi}$, qui en fait un torseur sous $\Lambda_n^{*}$ via (46).

(48) \struck{Le groupe} Essayons à prouver que les conditions
\marginal{en bas à gauche, en regard de (48) : Ceci \uncertain{vaut} aussi \uncertain{(48)} i.e.\ \uncertain{pour} le cas de l'octaèdre !}

\page{146}
d'orientation \struck{\ill{}} \uncertain{sg$_i$} \uncertain{$(i \in \pi)$} sur $G_0$ sont \uncertain{les mêmes} (cf plus haut). Si $n$ impair, cela \uncertain{provient} de ce que, sur $G_0$, engendré par \uncertain{des} \struck{\ill{}} générateurs $\tau_0, \tau_1, \tau_\infty$ satisfaisant aux relations (40) (41), il y a \struck{\ill{}} \add{\ill{}} un unique sous-groupe d'indice 2, de façon plus précise
\[
(49) \qquad (G_0)_{\mathrm{ab}} \simeq \text{\struck{\ill{}}}\ \pm 1 \qquad \text{si } n \text{ impair}
\]
[[\emph{NB} Dans tous les cas, $(G_0)_{\mathrm{ab}}$ est \struck{\ill{}} le groupe \add{engendré} par \uncertain{les} trois él.\ d'ordre 2, $\dot\tau_0, \dot\tau_1, \dot\tau_\infty$ (et $\dot\rho_\infty^{3} = 1$ \struck{\ill{}} implique $\dot\rho_\infty = 1$ i.e.\ $\dot\tau_1 \dot\tau_0 = 1$ i.e.\ $\dot\tau_0 = \dot\tau_1$, donc $(G_0)_{\mathrm{ab}}$ est engendré par $\dot\tau_1, \dot\tau_\infty$, avec la relation
\[
(50) \qquad (\dot\tau_1 \dot\tau_\infty)^{n} = 1, \quad \text{i.e.} \quad \dot\tau_1 = \dot\tau_\infty \ \text{si } n \text{ impair}
\]
ce qui montre (49) pour $n$ impair.

Pour $n$ pair, on a
\[
(51) \qquad (G_0)_{\mathrm{ab}} \simeq \{\dot\tau_1, \dot\tau_\infty \mid \dot\tau_1^{2} = \dot\tau_\infty^{2} = \uncertain{\dot\tau_1 \dot\tau_\infty}\} \simeq \mathbb{F}_2 \times \mathbb{F}_2 ,
\]
\note{la relation entre accolades est transcrite telle qu'elle se lit ; on attend $\dot\tau_1^2 = \dot\tau_\infty^2 = 1$}
il y a donc \add{un ens.\ $X$ de} \emph{trois} \uncertain{caractères} \struck{\ill{}} non triviaux $G_0 \to \pm 1$, qui \struck{\ill{}} \add{a priori} \uncertain{pourraient être permutés} \add{par $G$}, opérant sur $G_{0\,\mathrm{ab}}$ via $G/G_0 \simeq \struck{\ill{}}\ Z'_n$. Mais comme $Z \subset G$ centralise $G_0$, $\mathrm{Im}(Z \to Z'_n)$ \uncertain{($\simeq Z'^{\,2}_n$)} (par $\lambda \mapsto \lambda^2$) donc $Z'^{\,2}_n$ opère sur $X$ \uncertain{trivialement}, donc l'image dans $\mathfrak{S}_X \simeq \mathfrak{S}_3$ \uncertain{(de $Z'_n/Z'^{\,2}_n$)} \struck{\ill{}} est donc un 2-groupe \struck{\ill{} on voudrait prouver} \add{donc ou bien réduit} \add{à l'élément unité} (et \ill{}), ou bien : l'élément unité (et \ill{}) \struck{\ill{}} une transposition \struck{et on voudrait que} \add{Dans ce dernier cas} on \uncertain{aurait}

\page{147}\note{p. 75 de l'auteur.}
que ou bien les $\chi_i$ sont égaux — et égaux \add{alors} au \add{l'unique} \uncertain{caractère} de $X$ laissé fixe par l'action précédente — ou bien $\chi_i$ prend deux valeurs $\chi, \chi'$, correspondant à \struck{\ill{}} la partition de $\pi$ en deux classes suivant le \struck{\ill{}} ss-groupe de $Z'_n \simeq Z$ noyau du caractère
\[
\varepsilon : Z'_n \longrightarrow \{\pm 1\}
\]
qui décrit l'action de $Z'_n$ sur $X$. Mais il faudrait déjà se rendre compte s'il peut bien arriver que l'action de $Z'_n$ sur $X$ soit non triviale — dans le cas \uncertain{\ill{}} défini par \ill{} \ill{} de rg 2 $M$ sur $\Lambda_n$, \ill{} question se pose \ill{} pour le premier cas où la question se pose, savoir $n = 8$\note{« 8 » récrit sur « 5 »}, $Z'_n \simeq \{\pm 1\}$, quelle est la structure de $[\mathrm{Gl}'(2, \mathbb{Z}/8\mathbb{Z})^{\pm 1}]_{\mathrm{ab}}$ ? Tout ceci fait, il restera par ailleurs à \ill{} quelles \ill{} \ill{} \ill{} \ill{} il limite sur $n$ et \ill{} \add{pour la structure} \struck{l'action de relations} d'extension
\[
(52) \qquad 1 \longrightarrow \mathrm{Gl}'(2, \hat{\mathbb{Z}})^{\pm} \longrightarrow \mathrm{Gl}'(2, \hat{\mathbb{Z}}) \longrightarrow \hat{\mathbb{Z}}^{*}/\pm 1 \longrightarrow 1
\]
l'\struck{opération \ill{} d'extension} \add{\ill{} de} $\hat{\mathbb{Z}}^{*}/\{\pm 1\}$ sur $[\mathrm{Gl}'(2, \hat{\mathbb{Z}})^{\pm}]_{\mathrm{ab}} \simeq \mathbb{F}_2 \times \mathbb{F}_2$ (base $\dot\tau_1, \dot\tau_\infty$)
\[
(53) \qquad \begin{cases} \text{opération de } \hat{\mathbb{Z}}^{*} \ \ldots \\ \text{i.e.\ l'homomorphisme} \\ \hat{\mathbb{Z}}^{*}/\{\pm 1\} \longrightarrow \mathfrak{S}_{\{\dot\tau_1, \dot\tau_\infty, \dot\tau_1\dot\tau_\infty\}} \simeq \mathfrak{S}_3 \end{cases}
\]
\note{sous $\dot\tau_1 \dot\tau_\infty$, une ligne biffée et la correction « $\dot\lambda_0 = \dot\varepsilon_0$ », lecture douteuse}
il est clair, en regardant le sous-groupe de $\mathrm{G}(2, \hat{\mathbb{Z}})$ formé des matrices
\[
\Bigl\{ \begin{pmatrix} \lambda & \mu \\ 0 & 1 \end{pmatrix} \Bigm| \lambda \in \hat{\mathbb{Z}}^{*},\ \mu \in \hat{\mathbb{Z}} \Bigr\},
\]
qui est produit semi-direct de $\hat{\mathbb{Z}}^{*}$ (isomorphe au groupe des $\begin{pmatrix} \lambda & 0 \\ 0 & 1 \end{pmatrix}$)

\page{148}
et $\hat{\mathbb{Z}} . \dot\varepsilon_0$ (isomorphe au groupe des $\begin{pmatrix} 1 & \mu \\ 0 & 1 \end{pmatrix}$) que l'action de $\hat{\mathbb{Z}}^{*}$ sur $\dot\varepsilon_0$ est bien triviale — ce qui était clair a priori dans l'action de $\hat{\mathbb{Z}}^{*}/\pm 1$ dans (53) s'exprime par une \uncertain{caractère}
\[
(54) \qquad \varepsilon : \hat{\mathbb{Z}}^{*}/\pm 1 \longrightarrow \{\pm 1\}
\]
\note{sous $\{\pm 1\}$, un mot biffé, illisible}
qu'il faudrait expliciter, exprimant de quelle façon un élément arbitraire $\dot\lambda \in \hat{\mathbb{Z}}^{*}/\pm 1$ opère sur $\{\dot\tau_1, \dot\tau_\infty\}$ (avec $\dot\tau_1 = \dot\tau_0$) — rappelons qu'ici \ill{} \ill{} \ill{} \ill{} propres que
\note{le passage qui suit, depuis « rappelons » jusqu'au bas de la page, est barré de grands traits obliques ; plusieurs lignes en sont biffées en outre une à une}
\[
(55) \qquad r_0 = \begin{pmatrix} 0 & 1 \\ 1 & 0 \end{pmatrix}, \quad r_1 = \begin{pmatrix} -1 & 1 \\ 0 & 1 \end{pmatrix}, \quad r_\infty = \begin{pmatrix} -1 & 0 \\ 0 & 1 \end{pmatrix}.
\]
\struck{Si $\varepsilon$ n'est pas le \ill{} \ill{} trivial, \ill{} \ill{} $\hat{\mathbb{Z}}^{*}$ \ill{} \ill{}. La structure de \ill{} \ill{} \ill{}} \struck{Mise en \ill{}, les valeurs de $\varepsilon$ pour} \add{$\varepsilon_\ell$} \struck{\ill{}} \add{À vrai dire, il est clair a priori que} \struck{les valeurs} \struck{$\mathbb{Z}/\ell\mathbb{Z}$} \add{$\mathbb{Z}_\ell^{*} \simeq \mathbb{Z}_\ell \times \mathbb{F}_\ell^{*}$} \add{\ill{}} \add{$\mathbb{Z}_\ell^{*}$ \ill{} $\ell$ \uncertain{nombre premier}} \add{impair} est trivial, sur \struck{$\mathbb{Z}_\ell$} les \uncertain{\ill{}} [\uncertain{car il est trivial} \ill{} \struck{\ill{} qui \ill{} il \ill{} \ill{} \ill{} grandes}] \struck{$\mathbb{Z}_\ell$, ce qui montre que} Écrivons
\[
\hat{\mathbb{Z}}^{*} \simeq \hat{\mathbb{Z}}_2^{*} \times \underbrace{\hat{\mathbb{Z}}^{*\prime}}_{\prod_{\ell \text{ premier impair}} \mathbb{Z}_\ell^{*}},
\]
il est clair que si $\varepsilon \neq 1$, alors $\varepsilon$ se factorise par le facteur $\hat{\mathbb{Z}}^{*\prime}$ \uncertain{\ill{}} d'une \uncertain{caractère} sur \struck{$(\mathbb{Z}/\ell\mathbb{Z})^{*}$ avec \ill{} \ill{}} \uncertain{(i.e.\ il se factorise par d'une caractère sur $(\mathbb{Z}/\ell\mathbb{Z})^{*}$ avec $\ell$ premier impair)}, \uncertain{dans} \uncertain{ce} \ill{} \uncertain{les composantes} \uncertain{des caractères suivant} $\hat{\mathbb{Z}}_2^{*}$
\note{la fin de la page, sous les ratures, ne se laisse lire que par fragments ; l'ordre des insertions est incertain}

\page{149}\note{p. 76 de l'auteur.}
\struck{est non triviale. Donc pour voir si}
Notons que l'on \uncertain{détermine} déjà la valeur « universelle » de $\bigl(\mathrm{Gl}'(2, \mathbb{Z}/n\mathbb{Z})^{\pm}\bigr)_{\mathrm{ab}}$, \add{savoir $\mathbb{F}_2 \times \mathbb{F}_2$,} \note{un indice douteux, peut-être $n$, suit $\mathrm{Gl}'$} pour $n = 4$, de telle façon qu'il est \uncertain{garanti} \ill{} évident : a priori que $\varepsilon$ se factorise par
\[
(56) \qquad \hat{\mathbb{Z}}^{*}/\pm 1 \longrightarrow \underbrace{(\mathbb{Z}/4\mathbb{Z})^{*}/\pm 1}_{Z'_4 = \{1\}} \longrightarrow \{\pm 1\},
\]
\note{le numéro (56) est surchargé ; sous l'accolade, une ligne noircie de ratures, puis « $Z'_4 = \{1\}$ »}
\struck{et les deux questions \ill{} si la caractère est triviale, i.e.\ si $(\mathbb{Z}/8\mathbb{Z})^{*}/\pm 1$ \uncertain{$\simeq \pm 1$} \uncertain{opère} trivialement, ou non, sur $\bigl(\mathrm{Gl}'(2, \mathbb{Z}/8\mathbb{Z})^{\pm}\bigr)_{\mathrm{ab}} \simeq \mathbb{F}_2 \times \mathbb{F}_2$. Notons que $\mathrm{Gl}'(2, \mathbb{Z}/4\mathbb{Z})^{\pm} \simeq \mathfrak{S}_4 \times \pm 1$.}
\note{le passage biffé est en outre barré d'un trait oblique}
ce qui montre que la caractère $\varepsilon$ est triviale ! ]]
\note{la double crochet ferme l'aparté ouvert en p.~146 par « [[ NB »}

\struck{Revenons} Notons que $G_0^{+} = \mathrm{Aut}^{+}(S, F_{i_0})$ est le seul sous-groupe d'indice 2 de $G_0$ qui contienne $\rho_0$ — où $\rho_0$ est l'opération de rotation autour d'un sommet \uncertain{$s_0$}, \uncertain{associée} \uncertain{à} un repère $(s_0, s_1, s_\infty)$ \struck{\ill{}} \ill{} \uncertain{(\emph{NB} $\dot\rho_0 = \dot\tau_1 \dot\tau_\infty$) — pour les mêmes raisons} \uncertain{qu'on} \struck{\ill{}} $G_0^{+}$ ; il en est de même de $G_1^{+} = \mathrm{Aut}^{+}(S, F_{i_1})$, pour $i_1 = \lambda . i_0$ ; il suffit de voir que l'on \uncertain{a}
\[
(55) \qquad \rho_0 \overset{\mathrm{déf}}{=} \rho_{s_0, A_0} \in G_1^{+} = \mathrm{Aut}^{+}(S, F_{i_1}),
\]
\note{il numérote ici (55) une seconde fois}
pour ceci considérons l'étoile
\[
E = \mathrm{Et}_0(s_0) = \mathrm{Et}_1(\lambda . s_0) ;
\]
$\rho_0$ applique $E$ dans lui-même, et est un automorphisme \uncertain{(i.e.\ une « rotation »)} de la structure polygonale $A_{0E} = A_0 \cap \mathcal{P}_2(E)$, \struck{\ill{}} on en conclut qu'il est aussi un automorphisme de la structure polygonale $A_{1E} = A_1 \cap \mathcal{P}_2(E)$, puisque $A_{1E} = A_{0E}^{(\lambda)}$ — donc

\page{150}
$\rho_0$, qui fixe $\lambda s_0$ (puisqu'il fixe $s_0$ et commute à $\lambda$) et induit un automorphisme de l'étoile de $\lambda s_0$ \struck{\ill{}} pour $P_1$, est un automorphisme \add{\ill{}} de $P_1$ i.e.\ est $\in G_1^{+}$, ce qui achève de prouver
\[
(56) \qquad \begin{array}{l} \text{Les conventions d'orientation \add{sur $G_0$} pour les structures} \\ \text{de polyèdres $P_i$ ($i \in \pi$) sont les mêmes, i.e.\ il existe un} \\ \text{(unique) sous-groupe $G_0^{+}$ de $G_0$, tel que l'on ait} \\ \qquad G_0^{+} = \mathrm{Aut}^{+}(P_i) \qquad \forall\, i \in \pi . \end{array}
\]
\struck{Cela} achève en même temps de prouver (48).

Considérons maintenant l'action (par transport de structure) de $G = \mathrm{Aut}(S, F)$ sur $\vec{\pi}$, il est immédiat que cette action \add{commute} \struck{\ill{}} à la structure de $\Lambda_n^{*}$-torseur sur $\vec{\pi}$, d'où on déduit un unique homomorphisme
\[
(57) \qquad G \xrightarrow{\ \chi\ } (\mathbb{Z}/n\mathbb{Z})^{*} \qquad \text{la « \emph{caractère cyclotomique} »}
\]
caractérisé par la propriété
\[
(58) \qquad \forall\, \underbrace{i \in \pi,\ \omega \in \mu(i)}_{\text{d'où } \alpha = (i, \omega) \in \vec{\pi}},\ g \in G, \text{ on a} \quad \underbrace{g\overbrace{(i, \omega)}^{\alpha \in \vec{\pi}}}_{\text{transport de structure}} = \chi(g) . \underbrace{\overbrace{(i, \omega)}^{\alpha}}_{\text{opération arithmétique}} .
\]
\note{sous « opération arithmétique », la glose : « d'exposant $\chi(g)$ sur la structure $P_i = (S, A_i, F_i)$ »}
Le caractère $\chi$ factorise la caractère
\[
\chi' : G \longrightarrow (\mathbb{Z}/n\mathbb{Z})^{*}/\pm 1 \struck{\ \ill{}\ } = Z'_n \simeq Z \qquad (\text{« \emph{caractère arithmétique} »}),
\]
i.e.\ on a commutativité dans

\page{151}\note{p. 77 de l'auteur.}
\begin{tikzcd}
G \arrow[r, "\chi"] \arrow[dr, "\chi'"'] & (\mathbb{Z}/n\mathbb{Z})^{*} \arrow[d, "\mathrm{can.}"] \\
& Z \simeq (\mathbb{Z}/n\mathbb{Z})^{*}/\pm 1
\end{tikzcd}

(59) On voit que le noyau de $\chi$ est $G_0^{+}$, de sorte que la suite exacte (17) peut se déduire d'une suite exacte

\begin{tikzcd}[column sep=small]
1 \arrow[r] & G_0^{+} \arrow[r] & G \arrow[r, "\chi"] & (\mathbb{Z}/n\mathbb{Z})^{*} \arrow[r] & 1 \\
& & Z \simeq (\mathbb{Z}/n\mathbb{Z})^{*}/\pm 1 \arrow[u, hook] \arrow[ur, "{\dot\lambda \mapsto \lambda^{2}}"'] & &
\end{tikzcd}

\note{le numéro (60) est porté à gauche de cette suite ; sous la flèche $\chi$ : « car.\ cyclotomique »}
\marginal{en diagonale, à gauche : si il y a dans cas octaèdre ($n = 4$) il faut remplacer $Z$ par $\{1\}$}
\emph{NB} \struck{\ill{}} \add{\ill{}} par
\[
(61) \qquad G_0 = \bigl\{ g \in G \bigm| \chi(g) \in \{\pm 1\} \bigr\} \qquad \bigl[ = \mathrm{Ker}\bigl(\chi' : G \to (\mathbb{Z}/n\mathbb{Z})^{*}/\pm 1\bigr) \bigr]
\]
\emph{NB} On lit sur (60) que
\[
(62) \qquad \underbrace{Z \cap G_0^{+}}_{= \mathrm{Cent}(G_0^{+})} \simeq \bigl({}_2(\mathbb{Z}/n\mathbb{Z})^{*}\bigr)/(\pm 1) \simeq \mathbb{F}_2^{\,\mathbb{P}(n)'}/\mathrm{diag}
\]
\marginal{en diagonale, à gauche de (62) : y compris (cas $n = 4$), \struck{\ill{}} \ill{} $Z \cap G_0^{+} = \{1\}$ ; plus bas, « $Z' \cap \ldots \{1, a\}$ », douteux}
$\mathbb{P}(n)$ est l'ensemble des nbs premiers \struck{impairs} qui divisent $n$, et $\mathbb{P}'(n)$ est égal à $\mathbb{P}(n)$ sauf si $2 \in \mathbb{P}(n)$ et si $2$ intervient \struck{$n \equiv 0\ (2)$, $n \not\equiv 0\ (4)$, i.e.\ si \ill{} \ill{}} avec un \struck{\ill{}} exposant $e_2 \neq 2$ — auquel cas on pose
\[
(63) \qquad \begin{cases} \mathbb{P}'(n) = \mathbb{P}(n) \setminus \{2\} & \text{si } e_2 = 1 \\ \mathbb{P}'(n) = \mathbb{P}(n) \amalg \{\text{un élément supplémentaire}\} & \text{si } e_2 \geqslant 3 \\ (\mathbb{P}'(n) = \mathbb{P}(n) & \text{si } e_2 \in \{0, 2\}) \end{cases}
\]
donc
\[
(64) \qquad Z \cap G_0^{+} \simeq \mathbb{F}_2^{\,\nu} \quad \text{avec} \quad \nu = \underbrace{\nu(n)}_{\substack{\text{nb de facteurs} \\ \text{premiers de } n}} - 1 + \varepsilon, \qquad \varepsilon = \begin{cases} 0 & \text{si } e_2 \in \{0, 2\} \\ -1 & \text{si } e_2 = 1 \\ +1 & \text{si } e_2 \geqslant 3 . \end{cases}
\]
\note{dans $\nu = \nu(n) - 1 + \varepsilon$, le « $-1$ » est récrit sur un autre signe ; le « 0 » du premier cas est noirci d'une rature}

\page{152}
Notons que
\[
\nu \geqslant \nu(n) - 2 ,
\]
ce qui implique qu'on ne peut avoir $\nu = 0$, i.e.\ $G_0^{+} \cap Z$ que dans les cas suivants \add{$\nu(n) = 1$ \uncertain{ou} $\nu(n) = 2$, $n = 2p$ …}
\[
(65) \qquad G_0^{+} \cap Z = 0 \iff \begin{cases} n = 2p^{e} & p \text{ premier impair} \\ n = p^{e} & p \text{ premier impair} \\ n = 4 & \text{(cas octaédral)} \end{cases} \qquad e \geqslant 1
\]
\note{l'exposant $e$ de « $2p^{e}$ » est ajouté d'une encre plus forte ; le « 4 » de la troisième ligne est récrit}
\struck{Mais on fait attention \uncertain{que}, dans les cas} Dans le cas $n = 2^{e} \ldots$ où $e \geqslant 3$, on trouve \struck{$G_0^{+} \cap Z \simeq \mathbb{F}_2$.}
\[
(65) \qquad {}_2 Z \simeq \mathbb{F}_2^{2}, \qquad G_0^{+} \cap Z \simeq \mathbb{F}_2 \qquad (\text{pour } n = 2^{e},\ e \geqslant 3).
\]
\note{il numérote (65) une seconde fois}
On va encore \struck{définir} \add{introduire} l'application canonique
\[
(66) \qquad \vec{F} \xrightarrow{\ \vec{\tau}_F\ } \text{\struck{\ill{}}}\ \vec{\pi}
\]
s'insérant dans le diagramme commutatif (on fait un carré cartésien)

\begin{tikzcd}
\vec{F} \arrow[r, "\vec{\tau}_F"] \arrow[d, "\mathrm{can}"'] & \vec{\pi} \arrow[d, "\mathrm{can}"] \\
F \arrow[r, "\tau_F"] & \pi
\end{tikzcd}

\note{le numéro (67) est porté à gauche de ce carré, qui est marqué « cart » en son milieu}
(cf (24) pour le « type » des faces \struck{\ill{} $\tau_F(f, \ldots)$}), en notant que si $f \in F$, si $i = \tau_F(f)$ est son type, il y a correspondance 1-1 entre l'ens.\ $\vec{F}_f$ des deux orientations de $f$, et l'ens.\ $\mu(i)$ des deux orientables de $P_i$ … des deux éléments de $\vec{\pi}$ au-dessus de $i$. On \uncertain{trouve} de même un diagramme commutatif cartésien

\begin{tikzcd}
\vec{A} \arrow[r, "\vec{\tau}_A"] \arrow[d, "\mathrm{can}"'] & \vec{\pi} \arrow[d, "\mathrm{can}"] \\
A \arrow[r, "\tau_A"] & \pi
\end{tikzcd}

\note{le numéro (68) est porté à gauche de ce carré}

\page{153}\note{p. 78 de l'auteur.}
avec la compatibilité entre (67), (68) :
\[
(69) \qquad \text{si } \vec{a} \in \vec{A} \text{ est incident (\emph{avec orientations induites}) à } \vec{f} \in \vec{F}, \text{ alors } \vec{\tau}_A(\vec{a}) = \vec{\tau}_F(\vec{f})
\]
\uncertain{dans} \uncertain{un} \ill{} \uncertain{commutativité} en d'autres termes, dans

\begin{tikzcd}[column sep=small]
& \mathrm{Rep}(S, F) \arrow[dl, "\mathrm{can}"'] \arrow[dr, "\mathrm{can}"] & \\
\vec{A} \arrow[dr, "\vec{\tau}_A"'] & & \vec{F} \arrow[dl, "\vec{\tau}_F"] \\
& \vec{\pi} &
\end{tikzcd}

\note{le numéro (70) est porté à gauche de ce diagramme}
Je voudrais maintenant dire un peu plus la situation des diagrammes (60). Fixons un $s_0 \in S$, et considérons son stabilisateur \struck{\ill{}} $G_{s_0}$ dans $G$, \uncertain{on a donc}, divisant
\[
1 \longrightarrow \underset{\substack{\wr\ \text{non canoniquement} \\ \mathbb{Z}/n\mathbb{Z}}}{G^{+}_{0 s_0}} \longrightarrow G_{s_0} \longrightarrow (\mathbb{Z}/n\mathbb{Z})^{*} ;
\]
je dis que $G_{s_0} \xrightarrow{\chi \mid G_{s_0}} (\mathbb{Z}/n\mathbb{Z})^{*}$ est épi — c'est \uncertain{formel}, à partir du fait que $G_0^{+}$ est transitif sur l'ensemble $S$, d'où une \uncertain{suite} d'extensions canonique :
\[
(71) \qquad 1 \longrightarrow \underset{\substack{\wr \\ \mathbb{Z}/n\mathbb{Z}\ \text{\emph{non} can.}}}{G^{+}_{0 s_0}} \longrightarrow G_{s_0} \longrightarrow (\mathbb{Z}/n\mathbb{Z})^{*} \longrightarrow 1 .
\]
\note{sous la suite, deux gloses : à gauche « On a aussi des \uncertain{notations} \uncertain{propres} $\mathcal{U}_{s_0}$ (« \uncertain{unipotents} ») », sous $G_{s_0}$ « \uncertain{maximal} \uncertain{de} \uncertain{la} \ill{} \uncertain{B} (« Borel ») » ; à droite « mais il \uncertain{serait} \uncertain{commode} \ill{} \ill{} \ill{} \ill{} … », lecture très douteuse}

\page{154}
On aimerait scinder cette extension. \struck{Notons} Posons
\[
(72) \qquad \mathrm{Et}_F(s_0) \overset{\mathrm{déf}}{=} \coprod_{i \in \pi} \mathrm{Et}_i(s_0) = \bigl\{ s \in S \bigm| \{s_0, s\} \in A_0 \bigr\} ,
\]
et notons
\[
(73) \qquad \text{\struck{\ill{}}}\ G_{s_0} \text{ est transitif sur } \mathrm{Et}_F(s_0)
\]
ce qui résulte encore formellement du fait que $G$ est transitif sur $\vec{A}$, \struck{et $G_0^{+}$ est} \struck{transitif sur} \uncertain{i.e.\ le} stabilisateur \add{dans $G_{s_0}$} d'un élément $t_0 \in \mathrm{Et}_F(s_0)$, i.e.\ le stabilisateur de $a = (s_0, t_0) \in \vec{A}$, est un sous-groupe d'ordre 2, contenu dans $G_0^{+}$ \struck{\ill{}}, dit « \uncertain{diédral} », égal à $\pm 1$. \uncertain{Son} élément non trivial n'est autre que la notation $\tau_\infty$, plus précisément $\tau_\infty^{a_0}$ où $a_0 = \{s_0, t_0\}$, égal à $\tau_\infty^{r_0}$ où $r_0$ est un repère de la forme
\[
r_0 = (s_0, t_0, u_0) .
\]
Notons que le groupe
\[
Z \simeq (\mathbb{Z}/n\mathbb{Z})^{*}
\]
opère \emph{librement} sur $S$ (puisqu'il opère par automorphismes d'une structure d'espace homogène), \struck{\ill{}} sont $D_t$ \struck{\ill{}} l'orbite \struck{\ill{}} d'un $t \in S$. Notons que $Z$ stabilise $\mathrm{Et}_F(s_0)$ (pour $\forall\, s_0 \in S$), \uncertain{et} \uncertain{de} $Z$ (\uncertain{en} « \uncertain{voyons} ») \add{\uncertain{qui} \uncertain{se} \uncertain{rencontre}} $\mathrm{Et}_F(s_0)$, \uncertain{et} des toutes orbites (qui \uncertain{rencontrent}) … en plus \uncertain{contenues} \uncertain{dedans}. Choisissons donc, un plus de $s_0 \in S$, un \uncertain{voisin} \struck{$\ill{}$}
\[
(74) \qquad D_0 \ \text{\struck{\ill{}}} \subset \mathrm{Et}_F(s_0), \qquad \underbrace{D_0}_{Z\text{-orbite d'un } t_0 \in \mathrm{Et}_F(s_0)}
\]
\note{la ligne qui précède (74) est d'une lecture très incertaine}

\page{155}\note{p. 79 de l'auteur.}
Soit
\[
(75) \qquad T = T_{D_0} = \text{stabilisateur de } D_0 \text{ dans } G_{s_0} = \bigl\{ g \in G_{s_0} \bigm| \underbrace{g(D_0) = D_0}_{\text{i.e.\ } g(t_0) \in D_0} \bigr\}
\]
\struck{Posons} On a donc
\[
G^{+}_{0 s_0} = \text{\struck{\ill{}}}\ \subset G_{s_0} = \text{\struck{\ill{}}}\ B' ; \qquad T \subset G_{s_0} ,
\]
\note{sous $G_{s_0}$, une flèche d'inclusion verticale part de $T$ ; les deux termes biffés sont noircis (on devine $\mathcal{U}_{s_0}$ et $B$)}
\struck{quel est $T \cap G_0^{+}$} je dis que
\[
(76) \qquad T \cap G_0^{+} = \{1\}, \qquad T \xrightarrow{\ \chi \mid T\ } (\mathbb{Z}/n\mathbb{Z})^{*} \ \text{épi donc iso}
\]
en d'autres termes, je dis que $T$ définit un \emph{scindage} de la structure d'ext.\ (71)
\[
(77) \qquad G_{s_0} \simeq \text{\struck{\ill{}}}\ \underbrace{\mathcal{U}}_{\substack{G^{+}_{0 s_0} \simeq \mathbb{Z}/n\mathbb{Z} \\ \text{non can.!}}} . \underset{\substack{\wr \\ (\mathbb{Z}/n\mathbb{Z})^{*}}}{T} \qquad \text{½ direct} .
\]
Soit donc $g \in G^{+}_{0 s_0}$ qui soit dans $T$ i.e.\ tel que $g(t_0) \in D_0$ donc $g(t_0) = \dot\lambda t_0$ ($\dot\lambda \in Z$), \struck{comme \ill{} $g \in G_0$ et $g$ fixe $s_0$, \ill{} \ill{} il s'ensuit} \struck{\ill{} $\{s_0, t_0\}$ et $\{g s_0 = s_0, g(t_0)\}$ \ill{}} \struck{\ill{} structure $P_{i_0}$ ($i_0 \in \pi$)}, mais $g(t_0) = \lambda t_0$ \struck{signifie} \add{implique} \add{\ill{}} \add{suit}
\[
i_0 = \tau_A(\{s_0, t_0\}), \qquad i_1 = \tau_A(\{\underset{\substack{\parallel \\ g(s_0)}}{s_0}, g(t_0)\})
\]
comme on a $g \in G_0$ et $g(s_0) = s_0$, \ill{} \ill{} \ill{}
\[
i_1 = g(i_0) = i_0
\]
et comme $g(t_0) = \lambda . t_0$, \ill{}
\[
i_1 = T_\lambda . i_0 \qquad \text{d'où} \quad \lambda = 1, \ \text{i.e.\ } g(t_0) = t_0
\]
\struck{[\ill{} \ill{} \ill{} \ill{} : deux points $A$-adjacents} \struck{\ill{} d'un voyage, \ill{} \ill{} \ill{} \ill{}] \ill{} \ill{} \ill{} \ill{}}
\drawing{en bas à droite, dans l'encadré biffé : un point $s$ relié à deux points $t_1$ et $t_0$ d'une droite $D$}
\note{les quatre dernières lignes, encadrées, sont biffées}

\page{156}
[On pourrait dire que

(78) Soit $s_0 \in S$, $D$ un voisin. Alors si un $t_0 \in D$ est $A$-adjacent à $s_0$, tous les $t \in D$ le sont. L'application
\[
D \longrightarrow \pi, \qquad t \longmapsto \tau_A(\{s_0, t\})
\]
est \emph{bijective}.]

Ainsi $g \in G_0$ fixe $s_0$ et $t_0$, qui sont adjacents, donc si on avait $g \neq \mathrm{id}$, on aurait $g = \tau_\infty^{\{s_0, t_0\}}$, or $\chi(\tau_\infty) = -1$ \struck{et} tandis que \struck{\ill{}} l'on a supposé $g \in G_0^{+}$ i.e.\ $\chi(g) = 1$, donc on doit avoir $g = 1$.

Prouvons enfin $T \xrightarrow{\ \chi \mid T\ } (\mathbb{Z}/n\mathbb{Z})^{*} \simeq Z$ surjectif. Soit $\dot\lambda \in (\mathbb{Z}/n\mathbb{Z})^{*}$ \add{et soit $\lambda$ son image dans $Z = (\mathbb{Z}/n\mathbb{Z})^{*}/(\pm 1)$,} alors $\exists\, g \in G_{s_0}$ tel que l'on ait
\[
\text{\struck{\ill{}}}\ \text{ou} \quad g t_0 = \dot\lambda . t_0
\]
puisque $G_{s_0}$ est transitif sur $\mathrm{Et}_F(s_0)$, et cela implique que $g \in T$ (75). \struck{\ill{}} La relation \struck{\ill{}} signifie aussi
\[
\chi'(g) = \dot\lambda \quad \text{i.e.} \quad \chi(g) \in \pm \lambda ,
\]
mais d'autre part il y a ambiguïté dans le choix de $g$, via la possibilité
\[
g \longmapsto \text{\struck{\ill{}}}\ g \circ \tau_\infty , \qquad \text{donc} \quad \chi(g) \longmapsto \chi(g)\chi(\tau_\infty)
\]
or
\[
\chi(\tau_\infty) = -1
\]
donc en corrigeant au besoin $g$, on gagne !
\note{dans ce paragraphe, les rôles de $\lambda$ et $\dot\lambda$ sont échangés par rapport à l'ajout interlinéaire ; transcrit tel quel}

\emph{Remarque.} Cette façon de trouver une décomposition de $G_{s_0}$ en produit ½ direct $G^{+}_{0 s_0} . \text{\struck{\ill{}}}\ T$, où $T \subset G_{s_0}$ est tel que $T \to \tilde{Z} = G/G_0^{+}$, est « formelle » à partir des seules conditions T~1 à T~4, et l'axiome de régularité et d'orientabilité (T~8),
\note{la phrase continue à la page suivante}

\page{157}\note{p. 80 de l'auteur.}
\marginal{en diagonale, à gauche, deux notes : « cette \uncertain{présentation} \uncertain{précise} \ill{} \ill{} \ill{} \ill{} \ill{} \ill{} structure \ill{} $P_0$ \ill{} \ill{} \ill{} \ill{} \ill{} condition de régularité » ; puis « cette condition implique que $Z$ \uncertain{est} \uncertain{commutatif} »}
: l'exclusion des axiomes « arithmétiques » T~5) T~6) T~7), permettant entre autres d'identifier $\tilde{Z}$ à $(\mathbb{Z}/n\mathbb{Z})^{*}$, et d'interpréter l'action de $Z \simeq Z'_n$ sur $(S, F)$ par voie « arithmétique ».

Dans cet esprit, partons de la situation suivante : un 2-polyèdre triangulé comme régulier $P_0 = (S, A_0, F_0)$, son groupe d'autom.\ $G_0$, et un \emph{groupe $Z$ donné d'avance}
\[
Z \subset \mathfrak{S}_S , \qquad Z \subset \underset{\mathfrak{S}_S}{\mathrm{Centr}}(G_0) ,
\]
satisfaisant à la condition supplémentaire (\uncertain{peut-être} plus \uncertain{forte}, \uncertain{mais} \uncertain{vérifiée} \uncertain{dans} \uncertain{les} \uncertain{cas} \uncertain{plus} \uncertain{haut}, \uncertain{plus} la condition \uncertain{de} \uncertain{commutation} \uncertain{de} \uncertain{régularité})
\[
(79) \qquad \forall\, \lambda \in Z,\ \exists\, f \in \mathfrak{S}_S \ \bigl(\underline{\text{centralisant}\ Z}\bigr) \text{ telle que } \forall\, s, t \in S, \text{ on ait}
\]
\[
\{s, t\} \in A_0 \iff \{f^{-1} s, f^{-1} \lambda t\} \in A_0
\]
(cf transcription pp~66' (8)).

Alors par les considérations générales maintes fois répétées, on construit, pour tout $\lambda \in Z$, un $A_\lambda \subset$ \struck{\ill{}} $\mathcal{P}_2(S)$, et une opération \add{$\lambda \mapsto T_\lambda$} de $Z$ sur l'ens.\ $\pi$ ($\subset \mathcal{P}(\mathcal{P}_2(S))$) des $\{A_\lambda\}$, qui fait de $\pi$ un torseur sous $Z$ (sauf dans le cas octaédral, où on a $\pi = \struck{\{1\}}$ \add{réduit à un point}, tandis que l'on peut avoir $Z = \{1, \varepsilon\} \neq 1$).

Si l'on pose
\[
(80) \qquad G = \bigl\{ g \in \mathfrak{S}_S \bigm| g \text{ centralise } Z, \text{ et applique } \pi \text{ dans lui-même} \bigr\}
\]
alors $G$ opère sur $\pi$ par $Z$-translations, et on
\note{« pp~66' (8) » : renvoi à sa propre pagination, p.~66' ; la phrase continue à la page suivante}

\page{158}
\marginal{en diagonale, à gauche : On va \uncertain{supposer} \uncertain{l'existence} de la suite \uncertain{exacte} \uncertain{sauf} \uncertain{dans} le cas octaédral, où $Z = \{1, \varepsilon\}$.}
trouve une structure d'extension

\begin{tikzcd}
1 \arrow[r] & G_0 \arrow[r] & G \arrow[r, "\chi'"] & Z \arrow[r] & 1 \\
& & Z \arrow[u, hook] \arrow[ur, "{\lambda \mapsto \lambda^{2}}"'] & &
\end{tikzcd}

\note{le numéro (81) est porté à gauche de cette suite}
On peut ensuite définir une notion d'\emph{équivalence} de deux structures de 2-polyèdres \struck{réguliers} triangulés \add{\uncertain{semblables}}, munis d'un groupe $Z$ \struck{\ill{}} de permutations de $S$ (\uncertain{jouant} le rôle d'« \emph{homothéties} ») — en particulier aussi de 2-polyèdres à opérateurs.
\note{la suite de la page est blanche}


\section{IX Graphes à homothéties}

\page{159}\note{p. 81 de l'auteur. Le titre est le sien : « IX Graphes à homothéties », « homothéties » écrit au-dessus de « opérateurs », biffé. Un nouveau paragraphe (IX) commence ici, avec une numérotation des formules qui repart de (1).}
\marginal{en diagonale, à gauche, en regard de (4) : \struck{Cont\ill{}} c) implique que si $\lambda, \mu \in Z$, alors $\forall\, s \in S$, on a $\mathrm{Et}_{A_0}((\lambda\mu) s) = \mathrm{Et}_{A_0}((\mu\lambda) s)$ \uncertain{i.e.} $\lambda\mu = \mu\lambda$ si on suppose \ill{} \ill{} \ill{} — \uncertain{En} \uncertain{tout} \uncertain{état} \uncertain{de} \uncertain{cause} \uncertain{on} \uncertain{suppose} $Z$ \uncertain{commutatif}, \ill{} \ill{} \ill{} \ill{} $S \to \mathcal{P}(S)$, $s \mapsto \mathrm{Et}_A(s)$ \ill{} \ill{} \ill{} \ill{} \ill{} \ill{} — lecture très douteuse}
Nous \struck{prendrons} appelons ici structure de \emph{graphe} sur un ens.\ de sommets $S$, la donnée d'une partie
\[
(1) \qquad A \subset \mathcal{P}_2(S)
\]
Nous appelerons \emph{graphe à \struck{opérateurs} \add{homothéties}} une structure sur un ens.\ de base $S$ (l'ens.\ de sommets) définie par la donnée de deux structures
\[
(2) \qquad A_0 \subset \mathcal{P}_2(S), \qquad \underset{\text{sous-groupe}}{Z \subset \mathfrak{S}_S}
\]
satisfaisant aux conditions qui suivent. Notons tout de suite qu'on \uncertain{suppose} \uncertain{nullement} que \struck{\ill{} $Z$ \ill{} \ill{} groupe de $\mathfrak{S}_S$} les \struck{\ill{}} $\lambda \in Z$ respectent la structure de graphe \struck{\ill{}} $A$ de $S$. On suppose par contre :
\[
(3)\ (\mathrm{GH}\,1) \qquad \forall\, (s, t) \in S,\ \lambda \in Z,\ \exists\, f \in \mathfrak{S}_S, \text{ centralisant } Z, \text{ tel que}
\]
\[
\{s, t\} \in A_0 \iff \{f^{-1} s, f^{-1} \lambda t\} \in A_0
\]
\note{le premier $f^{-1}$ est surchargé ; ordre des quantificateurs tel qu'il se lit}
Cette condition peut se décomposer en trois :
\[
(4)\ \mathrm{GH}\,1' \quad \begin{cases} \text{a) } \forall\, \lambda \in Z,\ \exists\, A_\lambda \in \mathcal{P}(\mathcal{P}_2(S)) \text{ structure de graphe sur } S \\ \quad \text{(évidemment unique) telle que } \forall\, s, t \in S, \\ \quad \{s, t\} \in A_0 \iff \{s, \lambda t\} \in A_\lambda \\ \text{b) Cette structure de graphe } A_\lambda \text{ est isomorphe à } A_0 . \\ \text{c) \ldots} \end{cases}
\]
\note{après a), une ligne biffée « [NB \ill{} \ill{} \ill{} de $(S, A_0)$ \ill{} \ill{} $(S, A_\lambda)$ \ill{} \ill{}] » ; c) est d'abord rédigé, puis biffé, sous la forme « Pour les \ill{}, $f \in \mathfrak{S}_S$ qui \uncertain{centralisent} \ill{} » et « On \ill{} \ill{} qu'il \uncertain{existe} un isom. $f$ de » ; la rédaction retenue suit}
c) \struck{isomorphisme} \ill{} la structure $(S, A_0, Z)$ \ill{} la structure $(S, A_\lambda, Z)$, et plus précisément un isom.\ \uncertain{induisant} l'identité sur $Z$, i.e.\ que \struck{\ill{}} \add{induise} par une $f \in \mathfrak{S}_S$ qui centralise $Z$.

\struck{On pose} On pose
\[
(5) \qquad \pi = \{ A_\lambda \mid \lambda \in Z \} \subset \mathcal{P}(\mathcal{P}_2(S)) , \qquad A_i \ \text{lieu de } A_\lambda .
\]
\note{sous (5), une ligne biffée : « \ill{} \ill{} \ill{} $i \in \pi$ \ill{} \ill{} \ill{} \ill{} »}

\page{160}
On voit alors que toutes les structures de graphe à homothéties $(S, A_i, Z)$, pour $i \in \struck{\ill{}}\ \pi$, satisfont la condition (GH~1), et on \uncertain{peut} \uncertain{les} \uncertain{décrire} en termes de \add{la seule donnée} ($\pi \subset \mathcal{P}(\mathcal{P}_2(S))$) une opération de $Z$ sur $\pi$, notée
\[
(6) \qquad \begin{cases} \underset{Z}{\lambda} \longmapsto \underset{\mathfrak{S}_\pi}{T_\lambda} & \text{définie par} \\ \forall\, s, t \in S,\ i \in \pi, \quad \{s, t\} \in A_i \iff \{s, \lambda t\} \in A_{T_\lambda i} . \end{cases}
\]
\note{le numéro (6) est récrit sur un autre ; un mot biffé avant « définie par »}
\struck{Soit \ill{} appelle \emph{multigraphes à opérateurs}} \add{Deux} \add{$(S, A_0, Z_0)$, $(S, A_1, Z_1)$} graphe\supplied{s} à homothéties \add{sur le même ens.\ de sommets} \struck{$G_0 = A_0$} \struck{existe $\lambda \in Z_0$ tel que l'on ait} \struck{s'ils peuvent se déduire l'un de l'autre par} \struck{$A_1 = T_\lambda(A_0)$, une opération} seront dits \emph{équivalents} si l'on a $Z_0 = Z_1$, soit $Z$, et s'il existe $\lambda \in Z$ tel qu'on ait $A_1 = T_\lambda A_0$, \struck{\ill{}} on encore si $A_1$ est $\in$ de l'ens.\ $\pi$ construit \add{par} $(S, A_0, Z)$. Il est précisément, en termes de relations \add{bien} \uncertain{connues} d'équivalence, plus immédiat que c'est \uncertain{une} restriction que la relation d'isomorphie — et \uncertain{aussi} que la relation d'isomorphie commutant à l'opération de $Z$. La donnée, \struck{\ill{}} sur un ens.\ $S$, d'une classe d'équivalence $\pi$ en un sens de graphes à homothéties, sera appelée un \emph{multigraphe à homothéties}. Posons
\note{la page s'arrête sur « Posons » et un « dans » isolé en bas ; l'argument se poursuit au-delà de ce lot}

\end{document}
